\section{Introduction}
Latent stochastic differential equations (SDEs) enable the unsupervised discovery of dynamical systems from noisy, irregularly sampled time series, a setting common across scientific domains. 
They posit an unobserved latent state that evolves continuously in time and generates observations at a finite set of time points \cite{archambeau2007gaussian, li2020scalable}. 
Latent SDEs have been applied to neuroscience \cite{elgazzar2024universal, hu2024modeling}, disease modeling \cite{durso2024probabilistic}, and reduced-order fluid dynamics \cite{course2023state}.
Fitting these models requires both learning the latent dynamics and approximating the posterior over continuous trajectories.

Variational inference (VI) translates the problem of fitting latent SDEs into an optimization problem. 
However, evaluating the variational objective (ELBO) requires computing an integral over the full time horizon (a path-space KL divergence).
Simulation-based VI algorithms estimate this term by sweeping across a time discretization at each gradient step, either by integrating an ODE to compute the posterior marginals or by simulating sample trajectories from the variational SDE.
A recent line of work, simulation-free VI  \cite{course2023state,course2023amortized,bartosh2025sde}, sidesteps sequential simulation using stochastic VI. 
Rather than integrating, they compute an unbiased estimate of the objective by sampling a random time together with a state at that time.
They make this estimate efficient by parameterizing the posterior through its one-time marginal distributions, i.e., the distribution of the latent state at each time, allowing states to be sampled directly without simulating the preceding trajectory.

In this work, we show that the family of approximate posterior distributions over which simulation-free VI optimizes is \emph{strictly smaller} than that available to simulation-based VI.
One-time marginals do not determine a path law.
Existing simulation-free methods choose a single compatible drift from many possibilities, implicitly restricting the variational posterior. 
When observations are sparse or noisy, this arbitrary choice can substantially degrade both posterior inference and learned dynamics.

We propose \emph{Helmholtz-SDE}, a simulation-free VI algorithm that removes this restriction. 
Starting from any drift compatible with the prescribed marginals, we compute a Helmholtz correction that preserves those marginals while optimizing the ELBO. 
The resulting posterior family remains simulation-free but spans a larger class of path laws, closing the gap to simulation-based VI.